\label{background}
\subsection{Hardware-Aware Neural Architecture Search (HW-NAS)}

\paragraph{Neural Architecture Search (NAS)}
Designing Deep Neural Networks (DNNs) by hand is cumbersome and brittle process, originally carried out by highly-specialized human experts, and later automatized via Neural Architecture Search (NAS).
Instead of relying on the combination of human ingenuity, heuristics, and trial-and-error to design DNNs, NAS provides an automated and systematic strategy to discover architectures optimizing a specified objective. 
In its seminal work,~\citet{zoph2016neural} demonstrated the feasibility of designing DNNs automatically via Reinforcement Learning (RL), using post-training performance metrics as reward signal.
While effective,~\citep{zoph2016neural}'s pioneering work most notably demanded thousands of GPU hours---a computational demand amply exceeding the capabilities of most institutions.
Subsequent research introduced more computationally efficient approaches, including weight-sharing paradigms~\citep{brock2017smash} and gradient-based methods~\citep{liu2018darts}.
A major goal of NAS is to reduce human intervention in selecting specialized architectures for novel tasks, data modalities, or deployment conditions. By systematically exploring a well-defined search space, NAS not only identifies high-accuracy networks but can also reveal unusual architectures that depart from human-engineered patterns. 
In the years following~\citet{zoph2016neural}'s work, NAS has expanded to address a variety of tasks within Deep Learning (DL), including language modeling, and sequence-to-sequence learning~\citep{elsken2019neural}. 
Despite these advances, the overall \emph{cost of search} remains a key concern, especially when architectures must be verified experimentally for downstream performance. In practice, it is common to search for architectures within cell-based searchspaces~\citep{zoph2018learning}, hereby reducing DNN design to identifying an optimal cell within a possibly very large searchspace of all possible combinations.
Under reasonable assumptions, the computational burden of NAS \(C \) can be modeled as the product of (i) the number of candidate architectures probed \( \#\text{networks}\), and (ii) the average \emph{cost to evaluate one network}, \( c \), yielding \( C = \sum_i^{\mathcal N} c_i \approx \netpool c \), for a pool of probed networks, \(\mathcal N \). 
The cost of many early methods~\citep{zoph2016neural,baker2016designing,real2017large} was mainly dominated by \( c \), as they required \emph{training} candidate architectures on a target task to measure performance. Several strategies have been introduced to make evaluation more tractable. 
Training a single hypernetwork, for instance, amortizes \(C \) by training a single overparameterized model from which many sub-architectures can be sampled~\citep{brock2017smash, pham2018efficient,cai2019once}, hereby reducing \(C \) via reducing \( c \). However, training such a hypernetwork is still prohibitively expensive in scenarios where computational resources are not abundant, and does not have access to thousands of GPU hours~\citep{cai2019once} when searching large spaces of candidate architectures.

\paragraph{Training-free (TF) metrics}
To tackle these limitations, the research community developed \emph{training-free metrics} (TF metrics)~\citep{mellor2021neural,jacot2018neural,chen2021neural,chen2023understanding}, which allow to drastically reduce the and cost and resources needed to evaluate candidate architectures by scoring them at initialization, without training.
Thus, TF metrics serve as proxies for post-training performance, guiding search toward promising architectures without requiring any update of any architecture parameters. In particular, in this work we focus on~\emph{NASWOT}~\citep{mellor2021neural}, which quantifies the expressivity of an untrained network by examining its weight initialization and the distinct activation patterns it can generate. By analyzing how the network divides input data---by measuring the number of "linear regions" formed by decision boundaries at initialization---NASWOT produces a one-shot score correlating with final performance.
To search large spaces of candidate architectures, \citet{cavagnero2023freerea} propose complemeting NASWOT with \emph{LogSynflow}~\citep{cavagnero2023freerea, tanaka2020pruning}, correlating the performance of a candidate design with the gradient flow of a loss function specifically designed to reflect the model's capacity.
Searching architectures using TF metrics substantially reduces the cost of probing each candidate architecture \( c \). In turn, TF metrics prove particularly useful to guide possibly sample-inefficient algorithms, such as RL~\citep{zoph2016neural}, allowing RL to rival more sample-efficient methods like evolutionary algorithms~\citep{real2017large} for cell-based NAS.
Cell-based NAS algorithms are usually compared on cell-based search spaces such as NAS-Bench-101~\citep{ying2019bench} or NATS-Bench~\citep{dong2021nats}.

\paragraph{Hardware-Aware NAS (HW-NAS)}
While traditional NAS typically focuses on metrics such as validation accuracy or model size, hardware-aware NAS (HW-NAS) seeks to incorporate hardware performance indicators characteristic of the deployment scenario of the model while designing DNNs~\citep{benmeziane2023efficient}. 
Work in HW-NAS typically considers optimizing post-compilation latency~\citep{wu2019fbnet}, memory footprint~\citep{king2025micronas} or energy consumption~\citep{speckhard2022neural}. In this work we exclusively focus on \emph{post-compilation latency}, as it is crucial when deploying DNNs to edge devices and resource-constrained environments where low latency inference is paramount, such as autonomous robotics and internet of things (IoT) devices. 
Taking into account hardware performance when designing DNNs introduces another dimension of complexity, as architectures that perform well on a given downstream task on the high-end devices used to train them---typically, high-end GPUs or even high-performance computing (HPC) infrastructures---might be infeasible due to unacceptable latency or resource overhead in different deployment scenarios. 
Nonetheless, designing architectures for real-world use does require taking into account the specifities of the deployment scenario, such as the device used, or the expected workload. HW-NAS seeks to address precisely this very challenge, complementing NAS' automated approach to DNN design with considerations for hardware constraints and performance metrics specific to the deployment environment~\citep{li2021hw}.
However, designing a comprehensive HW-NAS strategy concerned with downstream real-world performance of trained DNNs is a challenging endeavour. 
First, measuring hardware performance at inference-time typically involves low-level compiling and deploying candidate devices, which can prove challenging and for non-domain experts, other than being costly~\citep{li2021hw}.
Second, hardware performance profiles can vary significantly across devices due to diverse hardware pipelines, compilers, and operator libraries~\citep{li2021hw,benmeziane2023efficient}.
In practice, the same high-level logical operations---e.g., a 3x3 convolutional filter---are compiled very differently on different hardware platforms~\citep{li2021hw,benmeziane2023efficient,laube2022expect}. In turn, this makes it nontrivial to identify good architectures across diverse devices.
HW-NAS-Bench~\citep{li2021hw} was proposed to promote reproducibility and expedite research in hardware-aware design for multi-device setting. 
Similarly to NATS-Bench~\citep{dong2021nats}, HW-NAS-Bench provides a curated set of architectures along with measured latency and energy consumption across multiple real-world hardware targets, ranging from commercial edge devices (Raspberry Pi, Google Pixel3) to specialized accelerators (FPGA, Eyeriss).

\subparagraph{Measuring latency for HW-NAS}
Measuring actual inference latency often proves crucial in hardware-aware NAS, yet directly compiling and timing each candidate architecture can be a nontrivial and costly process~\citep{li2021hw}. In principle, obtaining a ground-truth latency measurement for a specific model-device pair entails performing low-level optimizations. Differences in compiler toolchains or library support exacerbate this burden, resulting in latency measurements that can vary across hardware generations or even across software updates.
As an alternative, many HW-NAS approaches rely on \emph{approximations} of post-compilation latency. 
A popular choice is using lookup tables (LUT), modelling architectures \( a = \{o_1, o_2, \dots \} \) as a set of fundamental operations, and aggregating the per-operation runtimes \( t_d(o_i), \ o \in \mathcal O \) to obtain the network's latency \( \ell_d(a) \) for a given device \( d \). 
Typically, operation-wise runtimes are typically aggregating as \(\ell (a) \approx \hat{\ell}(a) = \sum_{o \in a} t_d(o) \)~\citep{wu2019fbnet,cai2018proxylessnas}. 
LUTs are fast to use at query time, and cheap to obtain even for very large search spaces. 
For instance, while FBNet~\citep{wu2019fbnet} contains \(\sim10^{21} \) architectures, obtaining a LUT for FBNet only requires deploying its 9 fundamental operations.
However, LUT deliberately fail to capture interactions among operations within a given network, which results in poor correlation between approximated and true latency values~\citep{laube2022expect}.
A more flexible strategy is to learn a \emph{predictor model} for latency, where an architecture is encoded using a feature vector \( h(a) \), and the associated latency \( \ell(a) \) is regressed using supervised learning on pairs of \( \tau = \{h(a_i), \ell(a_i)\}_{i=1}^{K} \)~\citep{laube2022expect}. Once trained, the predictor can quickly infer latency for unseen architectures, thereby accelerating the NAS loop.
\citet{laube2022expect} report the Kendall-Tau correlation (KT) between true and approximated latency values of various predictors versus LUT varying the number of examples needed to train a regression model \( K \), and report any of the predictors used---ranging from gradient-boosted trees to MLPs---drastically outperforms LUT's KT correlation of 0.6 with as little as \(K=60\) examples, reaching KT values of 0.8 with \( K = 500 \). 
In turn, \citet{laube2022expect} show that \emph{(ii)} LUT introduce severe misestimations of true latency, with unsatisfactory correlation with real-world latency and \emph{(ii)} predictors trained on a few hundreads of examples outperform them significantly.
However, such predictive methods introduce approximation errors making them unreliable in risk-averse deployments such as autonomous robotics, aerospace, or industrial monitoring.
Further, both LUTs and predictors assume substantial upfront data collection---building the table or training the model requires a nontrivial set of measured (architecture, latency) samples. When the target hardware is scarce or each compilation is expensive, these preliminary steps may already pose a major overhead. Furthermore, if practitioners suspect that distribution shifts might occur (e.g., minor firmware updates), a learned or tabulated model may become out-of-date, again risking suboptimal selections.

In contrast, our approach circumvents the use of any learned predictor or LUT by constraining the total number of architectures probed to a small but carefully selected set, hereby limiting \( \netpool \) at test time. 
Because we only need to compile and measure a handful of candidate networks, we can reliably obtain high-fidelity latency evaluations for each one, for small \( C \). 
This avoids introducing approximation errors into the search loop and reduces the risk of committing to an inaccurately estimated architecture in mission-critical deployments.

\subsection{Reinforcement Learning (RL)}
Reinforcement learning (RL) is an area of machine learning concerned with training agents to act optimally in an environment that provides feedback via a reward signal~\citep{sutton1998reinforcement}. 
In general, RL agent act on environment's states \(s_t \in \mathcal S \), to taking an action \(a_t \in \mathcal A \). 
For each interaction \(t=1, \dots, T \) the agent receives a reward \(r_t \in \mathbb R\) reflecting the desirability of taking action \( a_t \) when faced with \( s_t \), in light of the agent's ultimate goal of maximizing the total sum of (discounted) rewards \(\sum_t^T \gamma^t r_t \).
Over time, RL agents attempt to learn a conditional probability distribution \(\pi: \mathcal S \times \mathcal A \mapsto [0, 1] \) ("policy", \(\pi(a \vert s)\)) following this objective of cumulative reward maximization.
Over discrete state and action spaces---typical in cell-based NAS where any \( s \in \mathcal S \) represents a given cell and actions \( a \in \mathcal A \) discrete modifications of candidate architectures---popular RL algorithms include trust-region approaches, such as Trust Region Policy Optimization (TRPO)~\citep{schulman2015trust} and Proximal Policy Optimization (PPO)~\citep{schulman2017proximal}.
Unlike REINFORCE~\citep{zoph2016neural}, trust-region approaches such as PPO iteratively learn \( \pi \) bounding the relative difference between successive polices. In this, trust-region approaches aim at avoiding too drastic approaches from which on-policy RL risks to not recover otherwise.

While RL-based NAS set an important precedent, it was quickly recognized that repeatedly training every candidate network rendered the search protocol too costly. Thus, researchers moved towards alternative solutions that attempt to reduce computational demand by limiting the number of networks probed \( \netpool \), such as evolutionary algorithms~\citep{real2017large,cavagnero2023freerea}.
While evolutionary algorithms (EA) have shown promise for single-device NAS~\citep{real2017large,cavagnero2023freerea}, they remain ill-suited to scenarios in which multiple target devices must be supported, especially when each device differs substantially in its architectural or compiler-level constraints. 
Standard EA-based NAS pipelines re-initialize or re-search from scratch whenever the hardware setting changes, thereby incurring overhead. Even when combined with hardware metrics, EAs is inherently ill-posed to transfer previously acquired knowledge to a new device without discarding substantial parts of the discovered solutions.

In contrast, RL has exhibited strong cross-domain adaptability in fields where cross domain adaptability is key, such as robotics~\citep{kober2013reinforcement,akkaya2019solving}, and we thus argue RL's capacity for \emph{transfer learning} is especially attractive in HW-NAS. 
Rather than fully re-searching architectures for each new deployment device, one can meta-train an RL controller on a distribution of possible hardware profiles and re-use this controller whenever a novel device or updated firmware is encountered. 
Crucially, this approach avoids the need for separate predictor models per device; the policy itself encodes adaptation strategies, sampling only a few real-world measurements to confirm the final architecture design.
In this, RL's cost is limited not by limiting \( c \), but rather \( \netpool \) at test-time.

\subsection{Related Works}

\paragraph{Once-for-All (OFA)}
A pioneering approach in low-cost HW-NAS Once-for-All (OFA)~\citep{cai2019once}. OFA amortizes NAS' costs \( C \) by training a single over-parameterized hypernetwork which is pruned via EA~\citep{real2017large} into "child" sub-networks, with each subnetwork inheriting weights from the super-network without requiring any additional training epochs~\citep{brock2017smash}.
In principle, this allows cheap evaluation of the candidate architectures at test-time (\( c \downarrow \implies C \downarrow \)), by amortizing the cost of training the hypernetwork---1.2k GPU hours on NVIDIA V100 GPUs---on multiple searches. OFA also allows for HW-NAS, using LUT to guide EAs in selecting architectures fitting specific computational constraints. 
However, OFA's success relies on the costly super-network pre-training phase, which requires access to GPU hours simply unavailable in most scenarios. Further, sub-network selection process is also subject to approximation errors in hardware metrics, consequent to using LUT to estimate latency, which can yield sub-optimal sub-network choices~\citep{laube2022expect}.

\paragraph{Hardware-adaptive Efficient Latent Predictor (HELP)}
Hardware-adaptive Efficient Latent Predictor (HELP)~\citep{lee2021hardware} proposes a similar approach of keeping \( C \) manageable, by learning a meta-predictor to estimate post-compilation latency for unseen architectures, thus focusing on \( c \downarrow \implies C \downarrow \) too. Rather than measuring the latency of every candidate architecture on the physical device, HELP trains a surrogate model to map architecture encodings to hardware costs. Then, by assessing as little as 10 architectures on the target device, HELP adapts the meta-learned predictor, thus severily limiting the \emph{per-probe cost} \(c\) of each architecture. 
Yet, constructing HELP's predictor entails collecting a non-trivial set of real-world measured samples \(\tau_d=\{h(a_i), \ell(a_i)\}_{i=1}^K\) across possibly many different devices, which can prove challenging for non-experts in hardware~\citep{li2021hw}. Further, if the deployment scenario---the hardware platform or software environment---changes, the predictor may require retraining or fine-tuning. Lastly, in risk-sensitive deployments even modest misestimation of latency can invalidate guarantees on latency bounds, hindering applicability of predictor-based approaches such as HELP. 

Both OFA and HELP aim to mitigate \( C \) by limiting the overhead associated with \(c \). OFA amortizes \( c \) by reusing super-network weights and LUT---at the expense of accuracy in latency estimation---while HELP uses a meta-learned latency regressor---at the expense of safety in latency estimation. 
By contrast, our method explicitly mitigates \( C \) by limiting the \emph{number} of high-fidelity probes \(\netpool \) at test-time, sidestepping super-network pre-training or predictive models altogether. 
This design choice inherently avoids approximation errors in hardware metrics, as each of the few final architectures is tested with real measurements. Hence, we shift complexity into a meta-trained RL controller (Section~\ref{sec:method}) that can generalize over device variations learning to design well-performing DNNs across various target platforms.